%==============================================================================
% Section 6: Conclusion
%==============================================================================
\section{Conclusion}
\label{sec:conclusion}

This paper presents MambaSOHPredictor, a unified Mamba-based architecture
evaluated under a cross-battery, cross-temperature protocol substantially more
demanding than the timestep-based splits common in the literature. The
long-sequence configuration ($L = 4096$) achieved an MAE of 1.52\,\% and an
RMSE of 1.97\,\% on the fully unseen, underrepresented \degC{4} cell B0048,
surpassing the 2\,\% industrial threshold and reducing the relative error by
69\,\% versus a CNN-LSTM baseline, while the lightweight window-30 configuration
(79k parameters) kept a sub-2\,\% MAE at 56.1\,ms mean latency (P95
$=$ 78.3\,ms) on a standard CPU.

The 16-fold leave-one-battery-out evaluation confirmed B0048 was not a
favourably biased choice, while revealing a clear failure mode---degraded
accuracy for cells confined to rarely observed low-SOH regimes---mitigated by
expanding training coverage rather than by architectural changes. For anomaly
detection, we transparently report a modest F1 of 0.34 on proxy labels,
positioning the module as a sensitive early-warning filter requiring human
verification rather than an autonomous classifier.

Overall, industrial-grade accuracy and real-time edge deployability are jointly
achievable within a single lightweight, pure-PyTorch framework under a rigorous
generalisation protocol. Future work will expand training coverage across
underrepresented temperature\,$\times$\,SOH combinations, add remaining useful
life (RUL) prediction, integrate SOH and anomaly outputs into a
maintenance-decision risk matrix, and develop a retrieval-augmented generation
(RAG) prescription layer -- extensions being pursued within the broader solar
battery monitoring system under development by the author group.
